\documentclass[10pt,a4paper]{amsart}
\usepackage[margin=19mm]{geometry}
\usepackage{amsmath,amssymb,mathtools}
\usepackage[scr=rsfs,cal=euler]{mathalfa}
\usepackage{xfrac}
\usepackage[unicode,hidelinks,bookmarks=false]{hyperref}
\newcommand{\ie}{i.e.\ }
\newcommand{\cf}{cf.\ }
\newcommand{\resp}{respectively\ }
\newcommand{\Subsection}{\subsection}
\providecommand{\nobreakdash}{\nobreak\hbox{-}}
\setlength{\emergencystretch}{2em}
\multlinegap0pt
\usepackage{amssymb}
\newtheorem{lem}{Lemma}
\newtheorem{thm}{Theorem}
\newcommand{\calF}{\mathcal{F}}
\newcommand{\tr}{\text{tr}} %trace \newcommand{\calF}{\mathcal{F}}
\title{A Semi-Classical Szeg\H{o}-type Limit Theorem for Toeplitz Operators}
\author{Trevor Camper, Mishko Mitkovski}
\date{Source: arXiv:2411.19298v4 (CC BY 4.0)}
\begin{document}
\maketitle

\noindent\emph{Source and licence.} A Semi-Classical Szeg\H{o}-type Limit Theorem for Toeplitz Operators. Version arXiv:2411.19298v4 (CC BY 4.0), distributed by arXiv under the Creative Commons Attribution 4.0 International licence, http://creativecommons.org/licenses/by/4.0/, as recorded in the arXiv licence metadata for this exact version and shown on the abstract page https://arxiv.org/abs/2411.19298v4. Full licence notice: \texttt{licenses/arxiv-2411.19298-CC-BY-4.0.txt}.\par\smallskip

\noindent\textit{Editorial scope.} Counted unit: the article's thm thm:general\_szego\_lc (source0.tex lines 2633--2654). The article's complete original proof of it is delivered with it (source0.tex lines 2656--2956, one \texttt{proof} environment of 301 lines). No mathematics is written, repaired or re-derived: the statement and the proof are the article's own lines, in the article's own notation, and no retained line is substituted or re-flowed.  Retained uncounted as the source-local context the proof needs: lines 1780--1792; lines 2423--2447.  Nothing was omitted from any retained span, so the package carries no omission record.  No retained line contains a citation, so the wrapper carries no bibliography and no bibliography entry.  No retained line contains a figure environment, an image-inclusion command or drawing code, so no image is delivered and the package declares no asset dependency.  Editorial formatting only, in addition: an AMS \texttt{amsart} preamble that restates the article's own declarations and macros used by the retained lines, namely the environments \texttt{lem}, \texttt{thm} on separate counters, together with the standard packages those lines need.  The retained environments print in this document as Lemma 1, Theorem 1 in order of appearance, whereas the article prints the same items with the numbering of its own full document.  The complete native source of the article as distributed in the arXiv e-print for this exact version is bundled unchanged as \texttt{sources/169306/source0.tex}.
\begin{lem}\label{lem:diff_convex}
Let $\psi$ be real-valued. 
\begin{itemize}
    \item[(i)] If $\psi\in C[a,b]$, then for every $\varepsilon>0$ there exist convex functions $\psi_1,\psi_2$ on $[a,b]$ such that 
    \[
        \|\psi-(\psi_1-\psi_2)\|_{\infty} < \varepsilon.
    \]
    \item[(ii)] If $\psi\in C_0[0,\infty)$, then for every $\varepsilon>0$ there exist convex functions $\psi_1,\psi_2\in C_0[0,\infty)$ such that
    \[
        \|\psi-(\psi_1-\psi_2)\|_\infty < \varepsilon.
    \]
\end{itemize}
\end{lem}

\begin{thm}\label{thm:convex_szego_lc} Let $\sigma,\eta\in L^1(X,\nu)\cap L^2(X,\nu)$ be real-valued with $\eta$ non-negative.
Assume that the Toeplitz operators $T_\sigma^\alpha$ and
$T_\eta^\alpha$ commute for all $\alpha$. 
\leavevmode
\begin{itemize}
\item[(i)] Suppose $\sigma \ge 0$ and $\psi \in C([0,\infty))$
is convex such that the limit
\[
\lim_{x\to\infty} \frac{\psi(x)}{x}
\]
exists. Then
\[
\lim_{\alpha\to\infty}
\frac{1}{c(\alpha)}
\tr\!\big(T_\eta^\alpha \psi(T_\sigma^\alpha)\big)
=
\int_X \eta(x)\psi(\sigma(x))\,d\nu(x).
\]

\item[(ii)] If $\sigma \in L^\infty(X,\nu)$ and
$\psi$ is convex and continuous on
$[\mathrm{ess\,inf}\,\sigma,\mathrm{ess\,sup}\,\sigma]$,
then the same limit holds.
\end{itemize}
\end{thm}

\begin{thm}\label{thm:general_szego_lc}
Let $\sigma,\eta\in L^1(X,\nu)\cap L^2(X,\nu)$ be real-valued with $\eta$ non-negative. Assume that the Toeplitz operators $T_\sigma^\alpha$ and $T_\eta^\alpha$ commute for all $\alpha$.

\begin{enumerate}
\item[(i)] If $\sigma$ is bounded and
$\psi$ is continuous on
$[\mathrm{ess\,inf}\,\sigma,\mathrm{ess\,sup}\,\sigma]$,
then
\[
\lim_{\alpha\to\infty}
\frac{1}{c(\alpha)}
\tr\!\big(T_\eta^\alpha \psi(T_\sigma^\alpha)\big)
=
\int_X \eta(x)\psi(\sigma(x))\, d\nu(x).
\]

\item[(ii)] If $\sigma\ge0$ and
$\psi\in C([0,\infty))$ such that the limit
$\lim_{x\to\infty}\psi(x)/x$ exists,
then the same limit holds.
\end{enumerate}
\end{thm}

\begin{proof}
We use Lemma~\ref{lem:diff_convex} together with
Theorem~\ref{thm:convex_szego_lc}.

\medskip
\noindent
\textbf{Case (i).}
Assume $\sigma$ is bounded and $\psi$ is continuous on
$[\mathrm{ess\,inf}\,\sigma,\mathrm{ess\,sup}\,\sigma]$.

By Lemma~\ref{lem:diff_convex}(i), for every $\varepsilon>0$
there exist convex functions $\psi_1,\psi_2$ on this interval such that
\[
\|\psi-(\psi_1-\psi_2)\|_\infty < \varepsilon.
\]
Set $\varphi := \psi-(\psi_1-\psi_2)$.
Then $\|\varphi\|_\infty < \varepsilon$.

By functional calculus,
\[
\psi(T_\sigma^\alpha)
=
(\psi_1-\psi_2)(T_\sigma^\alpha)
+
\varphi(T_\sigma^\alpha).
\]
Hence
\[
\frac{1}{c(\alpha)}
\tr\!\big(T_\eta^\alpha \psi(T_\sigma^\alpha)\big)
=
\frac{1}{c(\alpha)}
\tr\!\big(T_\eta^\alpha (\psi_1-\psi_2)(T_\sigma^\alpha)\big)
+
R_\alpha,
\]
where
\[
R_\alpha
=
\frac{1}{c(\alpha)}
\tr\!\big(T_\eta^\alpha \varphi(T_\sigma^\alpha)\big).
\]

To estimate $R_\alpha$ we use the trace inequality
\[
|\tr(AB)| \le \|B\|\, \tr(|A|),
\]
valid for trace-class $A$ and bounded $B$.
Since $\|\varphi(T_\sigma^\alpha)\|
\le \|\varphi\|_\infty < \varepsilon$,
we obtain
\[
|R_\alpha|
\le
\varepsilon \,
\frac{1}{c(\alpha)} \tr(|T_\eta^\alpha|).
\]

Since $\eta\in L^1(X,\nu)$, write
\[
\eta=\eta_+-\eta_-,
\qquad
\eta_+,\eta_-\geq0.
\]
Then
\[
\begin{aligned}
\|T_\eta^\alpha\|_{\mathfrak S_1}
&\leq
\|T_{\eta_+}^\alpha\|_{\mathfrak S_1}
+
\|T_{\eta_-}^\alpha\|_{\mathfrak S_1}\\
&=
\operatorname{Tr}(T_{\eta_+}^\alpha)
+
\operatorname{Tr}(T_{\eta_-}^\alpha)\\
&=
\int_X|\eta(x)|\,d\nu_\alpha(x)\\
&=
c(\alpha)\int_X|\eta(x)|\,d\nu(x).
\end{aligned}
\]
Consequently,
\[
|R_\alpha|
\leq
\varepsilon\,
\frac{\|T_\eta^\alpha\|_{\mathfrak S_1}}{c(\alpha)}
\leq
\varepsilon\int_X|\eta(x)|\,d\nu(x).
\]

On the other hand, Theorem~\ref{thm:convex_szego_lc} applied to the convex functions
$\psi_1$ and $\psi_2$ gives
\[
\frac{1}{c(\alpha)}
\operatorname{Tr}\!\left(
T_\eta^\alpha
(\psi_1-\psi_2)(T_\sigma^\alpha)
\right)
\longrightarrow
\int_X
\eta(x)(\psi_1-\psi_2)(\sigma(x))
\,d\nu(x).
\]

Since
\[
\left\|
\psi-(\psi_1-\psi_2)
\right\|_\infty<\varepsilon,
\]
we have
\[
\left|
\int_X
\eta(x)
\bigl((\psi_1-\psi_2)(\sigma(x))-\psi(\sigma(x))\bigr)
\,d\nu(x)
\right|
\leq
\varepsilon\int_X|\eta(x)|\,d\nu(x).
\]

Combining this estimate with the bound for $R_\alpha$, we obtain
\[
\begin{aligned}
&\limsup_{\alpha\to\infty}
\left|
\frac{1}{c(\alpha)}
\operatorname{Tr}\!\left(
T_\eta^\alpha\psi(T_\sigma^\alpha)
\right)
-
\int_X\eta(x)\psi(\sigma(x))\,d\nu(x)
\right|
\\
&\qquad\leq
2\varepsilon\int_X|\eta(x)|\,d\nu(x).
\end{aligned}
\]
Since $\varepsilon>0$ is arbitrary, it follows that
\[
\lim_{\alpha\to\infty}
\frac{1}{c(\alpha)}
\operatorname{Tr}\!\left(
T_\eta^\alpha\psi(T_\sigma^\alpha)
\right)
=
\int_X\eta(x)\psi(\sigma(x))\,d\nu(x).
\]


\medskip
\noindent
\textbf{Case (ii).}

Let
\[
\beta:=\lim_{x\to\infty}\frac{\psi(x)}{x},
\qquad
r(x):=\psi(x)-\beta x.
\]
Then
\[
\frac{r(x)}{x}\longrightarrow0
\qquad (x\to\infty).
\]

Fix $\varepsilon>0$. Choose $R>0$ such that
\[
|r(t)|\leq\varepsilon t
\qquad\text{for all }t\geq R.
\]
Let $\chi_R\in C_c([0,\infty))$ satisfy
\[
0\leq\chi_R\leq1,\qquad
\chi_R(t)=1\ \text{for }0\leq t\leq R,\qquad
\chi_R(t)=0\ \text{for }t\geq2R.
\]
Set
\[
r_R:=\chi_R r.
\]
Then $r_R\in C_0([0,\infty))$ and
\[
|r(t)-r_R(t)|\leq\varepsilon t
\qquad\text{for all }t\geq0.
\]

By Lemma~\ref{lem:diff_convex}(ii), for every $\delta>0$ there exist convex functions
$q_1,q_2\in C_0([0,\infty))$ such that, with
$q:=q_1-q_2$,
\[
\|r_R-q\|_\infty<\delta.
\]
Theorem~\ref{thm:convex_szego_lc} applied to $q_1$ and $q_2$ gives
\[
\frac{1}{c(\alpha)}
\operatorname{Tr}\!\left(
T_\eta^\alpha q(T_\sigma^\alpha)
\right)
\longrightarrow
\int_X\eta(x)q(\sigma(x))\,d\nu(x).
\]

Since $\eta\geq0$,
\[
\frac{1}{c(\alpha)}
\left|
\operatorname{Tr}\!\left(
T_\eta^\alpha(r_R-q)(T_\sigma^\alpha)
\right)
\right|
\leq
\delta\,
\frac{\operatorname{Tr}(T_\eta^\alpha)}{c(\alpha)}
=
\delta\int_X\eta\,d\nu,
\]
and
\[
\left|
\int_X\eta(x)(r_R-q)(\sigma(x))\,d\nu(x)
\right|
\leq
\delta\int_X\eta\,d\nu.
\]
Letting $\delta\to0$, we obtain
\[
\lim_{\alpha\to\infty}
\frac{1}{c(\alpha)}
\operatorname{Tr}\!\left(
T_\eta^\alpha r_R(T_\sigma^\alpha)
\right)
=
\int_X\eta(x)r_R(\sigma(x))\,d\nu(x).
\]

Because $T_\eta^\alpha$ and $T_\sigma^\alpha$ commute and
$T_\eta^\alpha\geq0$, they have a common orthonormal eigenbasis.
Using $|r-r_R|\leq\varepsilon\,\mathrm{id}$ on $[0,\infty)$, we obtain
\[
\left|
\operatorname{Tr}\!\left(
T_\eta^\alpha(r-r_R)(T_\sigma^\alpha)
\right)
\right|
\leq
\varepsilon\,
\operatorname{Tr}\!\left(
T_\eta^\alpha T_\sigma^\alpha
\right).
\]
Hence, by the linear case of Theorem~\ref{thm:convex_szego_lc},
\[
\limsup_{\alpha\to\infty}
\frac{1}{c(\alpha)}
\left|
\operatorname{Tr}\!\left(
T_\eta^\alpha(r-r_R)(T_\sigma^\alpha)
\right)
\right|
\leq
\varepsilon\int_X\eta(x)\sigma(x)\,d\nu(x).
\]
Similarly,
\[
\left|
\int_X\eta(x)(r-r_R)(\sigma(x))\,d\nu(x)
\right|
\leq
\varepsilon\int_X\eta(x)\sigma(x)\,d\nu(x).
\]
Letting first $\alpha\to\infty$, then $\delta\to0$, and finally
$\varepsilon\to0$, gives
\[
\lim_{\alpha\to\infty}
\frac{1}{c(\alpha)}
\operatorname{Tr}\!\left(
T_\eta^\alpha r(T_\sigma^\alpha)
\right)
=
\int_X\eta(x)r(\sigma(x))\,d\nu(x).
\]

Finally, applying the linear case to the term $\beta x$ yields
\[
\lim_{\alpha\to\infty}
\frac{1}{c(\alpha)}
\operatorname{Tr}\!\left(
T_\eta^\alpha\psi(T_\sigma^\alpha)
\right)
=
\int_X\eta(x)\psi(\sigma(x))\,d\nu(x).
\]

%%%%%%%%%%%%%%%%%%%%%%%%%%%%%

\end{proof}

\end{document}
